
Tidal debris in the Milky Way and M31 as probes of dark matter:
Dark matter models beyond CDM predict deviations in the power spectrum as well as non-gravitational interactions. These lead to a variety of downstream effects on the abundance, structure, and evolution of dark matter halos. Dark matter subhalos undergo further evolution in the form of tidal mass loss as they orbit in the gravitational potential of their host halos. The observed satellite galaxies hosted by these subhalos are necessarily on a trajectory of similar tidal evolution from intact satellites to disrupted stellar streams and phase-mixed components of the stellar halo. Alternative dark matter models that affect the structure and interactions of dark matter halos have the potential to strongly affect the tidal evolution of the subhalos and satellites. Self-interacting dark matter models can form cores and core-collapsed cusps that lead to strong variability in the potential tidal disruption experienced by different systems. Baryonic feedback that also cores subhalos would potentially depend differently on subhalo masses. Warm dark matter with less concentrated regions and less subhalos of lower masses would also impact the field of satellite tidal debris. Prompt cusps might further lead to less tidal disruption. We intend to model the effects that changing density profiles in different dark matter models have on the observed satellite tidal debris and stellar halo of the Milky Way and M31 with Spec-S5.

Our goal is to predict the disrupted stellar mass function, the progenitor stellar mass function of disrupted satellites, and the distribution of fractional stellar mass loss experienced by satellites. We are doing this using a flexible model for mass loss that is combined with cosmological simulations to model the impact of CDM and SIDM-like dark matter subhalo profiles on the mass loss of the subhalo populations in systems of Milky Way and M31 mass (This is roughly the stage I have reached for the CDM case). These predictions will further have uncertainties propagated as well as marginalize over uncertainties in host disk mass, growth history, and potentially baryonic feedback. Combining the CDM and modified dark matter mass loss using a base of a range of Milky Way mass zoom-in simulations with stellar tidal tracks will result in the previously described statistics that should differ between dark matter models. We then plan to match the functions of fractional stellar mass loss to star-tagged realizations of the same cosmological dark matter only simulations with different galaxy-halo connections. For example, matching the SIDM more extreme stellar mass loss with a puffy galaxy-halo connection that results in a similar distribution. We can then use the star-tagged realistic stellar halos produced with that galaxy-halo connection and combine with a mock generation framework that upsamples star particles to stars. We can determine metrics at this or the previous stage for the “disruptedness” of the stellar halo (correlation functions, other methods for filamentary, cohesive, and clumpy vs. diffuse structure). Mock observations of these systems can then constrain Spec-S5’s power in using these metrics on the 6D stellar halo and tidal debris (including the radial velocity phase space) of the Milky Way and M31 to constrain dark matter models.
